\documentclass[11pt,a4paper]{article}

\usepackage[utf8]{inputenc}
\usepackage[T1]{fontenc}
\usepackage[french]{babel}
\usepackage[margin=2.5cm]{geometry}
\usepackage{graphicx}
\usepackage{booktabs}
\usepackage[hidelinks]{hyperref}

\newcommand{\acompleter}[1]{\par\noindent\textit{À compléter : #1}\par}

% Les figures absentes sont remplacées par un cadre pour garder le document compilable.
\newcommand{\figureattendue}[3]{%
  \begin{figure}[htbp]
    \centering
    \IfFileExists{#1}{%
      \includegraphics[width=0.9\linewidth,height=0.42\textheight,keepaspectratio]{#1}%
    }{%
      \fbox{\parbox[c][3cm][c]{0.85\linewidth}{\centering
        Figure à exporter depuis le notebook.\\
        Fichier attendu : \texttt{\detokenize{#1}}}}%
    }
    \caption{#2}
    \label{#3}
  \end{figure}%
}

\title{IFT6758 --- Devoir 4\\Analyse de sentiments sur les critiques IMDb}
\author{Nom : \textbf{Marsel Bakashov}\\Numéro étudiant : \textbf{20216284}}
\date{\today}

\begin{document}
\maketitle
\noindent\textbf{Brouillon de travail --- analyses et résultats à compléter.}
\tableofcontents
\newpage

\section*{T1.3 --- Préparation des données et exploration initiale}
\addcontentsline{toc}{section}{T1.3 --- Préparation des données et exploration initiale}

\begin{table}[htbp]
\centering
\begin{tabular}{ll}
\toprule
Propriété & Résultat \\
\midrule
Nombre initial d'observations & 50\,000 \\
Nombre de valeurs distinctes de la cible & 2 : \texttt{positive}, \texttt{negative} \\
Valeurs manquantes par colonne & 0 dans \texttt{review} et \texttt{sentiment} \\
Lignes en double au-delà de la première occurrence & 418 \\
Nombre d'observations après déduplication & 49\,582 \\
\bottomrule
\end{tabular}
\caption{Résumé du jeu de données avant et après déduplication.}
\end{table}

Le jeu de données initial contient 50\,000 observations et deux colonnes :
\texttt{review} et \texttt{sentiment}. La cible comporte deux classes,
\texttt{positive} et \texttt{negative}, avec 25\,000 observations chacune
avant déduplication. La vérification avec \texttt{isna().sum()} ne relève
aucune valeur manquante dans les deux colonnes.

La méthode \texttt{duplicated()} identifie 418 lignes en double au-delà de
leur première occurrence, en considérant l'ensemble des colonnes.
La vérification sur la seule colonne \texttt{review} retrouve également
418 doublons. Après suppression des lignes dupliquées avec
\texttt{drop\_duplicates}, avec le paramètre \texttt{keep="first"}, le jeu de données contient
49\,582 observations. Les contrôles finaux ne détectent plus ni lignes
ni critiques dupliquées : les critiques conservées sont donc uniques.

\section*{T3.1 --- Répartition des sentiments}
\addcontentsline{toc}{section}{T3.1 --- Répartition des sentiments}
\figureattendue{figures/sentiments.png}{Effectifs des sentiments après déduplication.}{fig:sentiments}
Après déduplication, le jeu de données contient 24\,884 critiques positives
(50,19\,\%) et 24\,698 critiques négatives (49,81\,\%), sur un total de
49\,582 critiques. La figure~\ref{fig:sentiments} montre des effectifs très
proches : l'écart est de 186 critiques, soit environ 0,38 point de pourcentage
entre les proportions des deux classes. Le jeu de données reste donc
approximativement équilibré ; aucune classe ne domine fortement l'autre.
Cet équilibre ne suffit toutefois pas à établir que les critiques seront
faciles à classifier.

\section*{T3.2 --- Longueur des critiques}
\addcontentsline{toc}{section}{T3.2 --- Longueur des critiques}
\figureattendue{figures/longueurs.png}{Distribution du nombre de mots par critique prétraitée.}{fig:longueurs}
Les critiques prétraitées n'ont pas toutes la même longueur. Sur les
49\,582 critiques, le nombre de mots varie de 3 à 1\,435, avec une moyenne
de 119,67 mots et un écart-type de 90,35 mots. La médiane est de 89 mots ;
les premier et troisième quartiles sont respectivement de 64 et 145 mots.
La moitié centrale des critiques contient donc entre 64 et 145 mots.

L'histogramme de la figure~\ref{fig:longueurs} présente une asymétrie à
droite : la majorité des critiques est relativement courte, tandis que
quelques critiques très longues étirent la distribution et contribuent à
une moyenne supérieure à la médiane. Une longueur unique ne décrit donc
pas toute la variabilité du corpus.

Ces longueurs correspondent au nombre de mots séparés par des espaces
après prétraitement, et non au nombre de caractères ni aux tokens du
tokeniseur DistilBERT. La longueur maximale d'entrée du transformeur devra
être choisie à partir de sa propre tokenisation, en tenant compte du
compromis entre troncature des critiques et coût de calcul.

\section*{T3.3 --- Fréquence des mots et difficulté de classification}
\addcontentsline{toc}{section}{T3.3 --- Fréquence des mots et difficulté de classification}
\figureattendue{figures/mots_frequents.png}{Les 20 mots les plus fréquents après prétraitement.}{fig:frequents}
\figureattendue{figures/mots_rares.png}{Les 20 mots les moins fréquents après prétraitement.}{fig:rares}
Le corpus prétraité contient 92\,007 mots distincts. Les 20 mots les plus
fréquents, représentés dans la figure~\ref{fig:frequents}, sont, par ordre
décroissant : \texttt{movie}, \texttt{film}, \texttt{one}, \texttt{like},
\texttt{time}, \texttt{good}, \texttt{character}, \texttt{story},
\texttt{even}, \texttt{get}, \texttt{would}, \texttt{make}, \texttt{see},
\texttt{really}, \texttt{scene}, \texttt{well}, \texttt{much},
\texttt{bad}, \texttt{people} et \texttt{great}.
Les termes dominants, notamment \texttt{movie} et \texttt{film}, reflètent
le thème cinématographique du corpus, sans exprimer nécessairement un
sentiment. Des mots évaluatifs comme \texttt{good}, \texttt{bad} et
\texttt{great} peuvent en revanche fournir des indices utiles pour la
classification.

Les 20 mots affichés dans la figure~\ref{fig:rares} n'apparaissent chacun
qu'une seule fois. On y trouve notamment \texttt{seneca}, \texttt{skynet},
\texttt{eucharist} et \texttt{dinasours}, ce qui illustre la présence de
références particulières, de termes spécialisés ou de formes possiblement
mal orthographiées. Ils constituent une sélection parmi les mots de
fréquence minimale : le code retient les 20 dernières entrées de la série
triée par fréquence décroissante avec \texttt{value\_counts()}, sans
imposer de critère secondaire pour départager les ex æquo. Il ne s'agit
donc pas nécessairement des seuls mots de fréquence égale à un.

L'EDA suggère que la classification est envisageable, grâce à un corpus
approximativement équilibré et à la présence de vocabulaire évaluatif,
mais qu'elle ne sera pas triviale. Les mots de thème sont partagés entre
les sentiments, les longueurs varient fortement et les mots rares offrent
peu d'exemples pour apprendre une association fiable. La polarité peut
aussi dépendre du contexte, des négations ou de l'ironie, que des fréquences
globales ne permettent pas d'analyser. En particulier, la suppression des
stopwords peut retirer des négations et faire perdre une information
utile. Ces observations ne suffisent donc pas à prédire une accuracy :
les performances devront être mesurées sur l'ensemble de test.

\section*{T4.4 --- Représentation TF-IDF et explicabilité LIME}
\addcontentsline{toc}{section}{T4.4 --- Représentation TF-IDF et explicabilité LIME}
\subsection*{Limites de TF-IDF et représentation alternative}
Avec les paramètres par défaut de \texttt{TfidfVectorizer}, chaque critique
est représentée par les poids de ses mots individuels. Cette approche
ignore leur ordre et les expressions composées : elle ne représente donc
pas directement les relations syntaxiques ou la portée d'une négation.
Elle ne rapproche pas non plus automatiquement des synonymes et attribue
la même caractéristique à un mot quel que soit son contexte. Enfin,
l'espace de caractéristiques est de grande dimension et très creux,
car chaque critique ne contient qu'une petite partie du vocabulaire ; les
mots absents du vocabulaire appris à l'entraînement sont ignorés lors de
la transformation du test~\cite{sklearn-text}.

Une alternative est une représentation dense et contextualisée issue d'un
transformeur préentraîné de type BERT. Contrairement aux poids TF-IDF,
les représentations des tokens dépendent du contexte gauche et droit et
peuvent être adaptées à la classification par ajustement du
modèle~\cite{bert}. Cette approche pourrait mieux exploiter les relations
entre mots et certaines formulations complexes du sentiment, mais son
coût de calcul est plus élevé et une amélioration sur ce corpus doit être
vérifiée expérimentalement. La tâche 5 étudiera cette piste avec
DistilBERT. Aucun modèle ne peut toutefois reconstituer systématiquement
une négation supprimée lors du prétraitement : cette perte d'information
reste une limite commune aux données utilisées ici.

\subsection*{Interprétation d'une prédiction avec LIME}
\figureattendue{figures/lime.png}{Explication LIME de la critique sélectionnée dans le notebook.}{fig:lime}
La critique sélectionnée porte l'identifiant 29\,171 et sa classe réelle
est négative. Le modèle attribue à la classe positive une probabilité
arrondie de 0,0602, soit 6,02\,\%, contre environ 93,98\,\% à la classe
négative. La prédiction est donc négative et concorde avec l'étiquette
réelle. Il s'agit des probabilités estimées par le modèle, et non d'une
garantie de justesse ou d'une mesure démontrée de calibration.

LIME explique ici la sortie de la classe \texttt{positive}. Dans la
figure~\ref{fig:lime}, les contributions rouges, de signe négatif,
défavorisent cette classe ; les contributions vertes, de signe positif,
la favorisent. Parmi les mots affichés, \texttt{incoherent},
\texttt{existent}, \texttt{moronic}, \texttt{dreadful} et
\texttt{horrible} ont les contributions négatives les plus importantes.
Les termes \texttt{pawn}, \texttt{highly} et \texttt{waterproof}
contribuent dans le sens positif, mais leurs poids locaux sont moins
élevés en valeur absolue que ceux des principaux termes négatifs.
Plusieurs mots explicitement péjoratifs fournissent ainsi une explication
plausible de la prédiction négative. Le signe associé à un mot tel que
\texttt{existent} ne signifie toutefois pas qu'il serait intrinsèquement
négatif : il décrit une association locale apprise par le modèle.

LIME construit une approximation interprétable dans le voisinage de cette
critique, à partir de perturbations du texte~\cite{lime}. Les barres sont
les coefficients de cette approximation locale, pas les paramètres
originaux de Naïve Bayes, et leur somme ne doit pas être assimilée à la
probabilité prédite. L'explication dépend des perturbations et du
voisinage retenus ; elle ne constitue ni une explication causale ni une
description complète du comportement du classificateur sur tout le corpus.

\section*{T5 --- Exploration de DistilBERT et comparaison}
\addcontentsline{toc}{section}{T5 --- Exploration de DistilBERT et comparaison}
\subsection*{T5.1 --- Implémentation et protocole expérimental}
\acompleter{Décrire le modèle \texttt{distilbert-base-uncased}, les données prétraitées utilisées et l'encodage négatif = 0, positif = 1.}
\acompleter{Confirmer le même découpage entraînement/test que Naïve Bayes (\texttt{test\_size=0.2}, \texttt{random\_state=42}). Décrire la validation réservée dans l'ensemble d'entraînement pour choisir les hyperparamètres sans utiliser le test.}
\acompleter{Documenter la tokenisation, la troncature, le padding, la longueur maximale, la taille des lots, les graines aléatoires, les versions principales, le matériel et le temps d'entraînement mesuré.}

\subsection*{T5.2 --- Deux approches d'ajustement des hyperparamètres}
\acompleter{Décrire une configuration de référence, puis deux approches d'ajustement justifiées. Expliquer les hypothèses, les valeurs explorées, le critère de sélection sur validation et les résultats observés.}

\begin{table}[htbp]
\centering
\begin{tabular}{llll}
\toprule
Expérience & Configuration & Métrique de validation & Résultat \\
\midrule
Référence & À compléter & À compléter & --- \\
Approche 1 & À compléter & À compléter & --- \\
Approche 2 & À compléter & À compléter & --- \\
\bottomrule
\end{tabular}
\caption{Synthèse des expériences. Ajouter des lignes pour les configurations réellement évaluées.}
\end{table}
\acompleter{Justifier le choix du meilleur modèle à partir des résultats de validation.}

\subsection*{T5.3 --- Évaluation et comparaison des modèles}
\acompleter{Définir l'accuracy, la précision, le rappel et le score F1. Préciser la classe positive et la convention d'agrégation employée pour comparer les modèles de façon cohérente.}

\begin{table}[htbp]
\centering
\begin{tabular}{lcccc}
\toprule
Modèle & Accuracy & Précision & Rappel & F1 \\
\midrule
TF-IDF + Naïve Bayes & --- & --- & --- & --- \\
DistilBERT sélectionné & --- & --- & --- & --- \\
\bottomrule
\end{tabular}
\caption{Performances mesurées sur le même ensemble de test.}
\end{table}

\figureattendue{figures/confusion_nb.png}{Matrice de confusion de Naïve Bayes sur l'ensemble de test.}{fig:confusion-nb}
\figureattendue{figures/confusion_distilbert.png}{Matrice de confusion du meilleur DistilBERT sur le même ensemble de test.}{fig:confusion-distilbert}
\acompleter{Identifier les axes et l'ordre des classes des matrices. Comparer les faux positifs, les faux négatifs et les métriques ; discuter de l'ampleur de l'amélioration réellement obtenue sans supposer que DistilBERT sera meilleur.}
\acompleter{Discuter du compromis entre performances, coût d'entraînement et interprétabilité. Distinguer une amélioration numérique d'une éventuelle significativité statistique, qui demanderait une analyse supplémentaire.}

\addcontentsline{toc}{section}{Références}
\begin{thebibliography}{9}
\bibitem{sklearn-text}
Scikit-learn, \emph{Feature extraction --- Text feature extraction},
sections sur les représentations en sac de mots, TF-IDF et leurs limites.
\url{https://scikit-learn.org/stable/modules/feature_extraction.html#text-feature-extraction}.

\bibitem{bert}
J. Devlin, M.-W. Chang, K. Lee et K. Toutanova (2019).
\emph{BERT: Pre-training of Deep Bidirectional Transformers for Language Understanding}.
NAACL-HLT, p.~4171--4186.
\url{https://aclanthology.org/N19-1423/}.

\bibitem{lime}
M. Ribeiro, S. Singh et C. Guestrin (2016).
\emph{``Why Should I Trust You?'': Explaining the Predictions of Any Classifier}.
NAACL-HLT: Demonstrations, p.~97--101.
\url{https://aclanthology.org/N16-3020/}.
\end{thebibliography}

\end{document}
